\toolchapter{G}{Stability of nonlinear systems}{Poles belong to linear models. When drag is quadratic, motors saturate and bodies rotate, stability needs a tool that does not solve the equation: a function that can only go down.}
\label{app:stab}

\lead{\usedcard{\setuprow{Lesson I.2}{stability defined; linearisation; a limit cycle}\setuprow{Lesson I.4}{energy as a Lyapunov function, the invariance principle}\setuprow{Lessons I.3, I.13, I.14}{time scales, what separation promises}\setuprow{Lesson I.6}{a loop that switches}\setuprow{Parts II, III}{Lyapunov designs (II.10, II.18, III.7, III.15–III.17)}}%
The drone is not linear. Its drag grows with the square of the speed, its motors have limits, its attitude is a rotation. The lessons deal with this in two ways: they linearise, and where that is not enough they argue with energy. This appendix gives both arguments their exact form — when a linearisation may be trusted, how a Lyapunov function proves stability and how to find one, what to do when it only \emph{almost} proves it, and how far from the equilibrium the proof reaches.}

\section{Equilibria and linearisation}

Write the system as $\dot{\vx} = f(\vx)$. A point $\vx^*$ with $f(\vx^*) = 0$ is an \term{equilibrium}: started there, the system stays there. A nonlinear system can have several equilibria, or none; a linear one, $\dot{\vx} = A\vx$ with invertible $A$, has exactly one.

Near an equilibrium a smooth $f$ is well approximated by its first-order Taylor expansion. With the deviation $\symbf z = \vx - \vx^*$,
\begin{equation}\label{eq:G-lin}
  \dot{\symbf z} = A\,\symbf z + (\text{terms of second order in } \symbf z), \qquad
  A = \frac{\partial f}{\partial\vx}\Big|_{\vx^*} = \begin{pmatrix}\dfrac{\partial f_1}{\partial x_1} & \cdots & \dfrac{\partial f_1}{\partial x_n}\\ \vdots & & \vdots\\ \dfrac{\partial f_n}{\partial x_1} & \cdots & \dfrac{\partial f_n}{\partial x_n}\end{pmatrix}_{\vx^*} .
\end{equation}
The matrix $A$ is the \term{Jacobian} of $f$, and $\dot{\symbf z} = A\symbf z$ is the \term{linearisation} at $\vx^*$. With an input, $\dot{\vx} = f(\vx, \symbf u)$, the same expansion about an equilibrium pair $(\vx^*, \symbf u^*)$ gives $\dot{\symbf z} = A\symbf z + B\symbf v$ with $B = \partial f/\partial\symbf u$ and $\symbf v = \symbf u - \symbf u^*$.

\Cref{thm:spring-indirect} says what the linearisation is worth: if all eigenvalues of $A$ have negative real parts, the equilibrium of the \emph{nonlinear} system is asymptotically stable; if one has a positive real part, it is unstable; if the rightmost ones sit on the imaginary axis, the linearisation decides nothing.

\begin{example}[Linearising the drag]
\label{ex:G-drag}
Vertically the drone obeys $m\dot v = T - mg - c_v|v|v$. Linearise it about a steady climb at the speed $v_0 > 0$, and about hover.

\textbf{Step 1 — the equilibrium.} A steady climb needs $T_0 = mg + c_vv_0^2$.

\textbf{Step 2 — the derivatives.} $f(v, T) = (T - mg - c_v|v|v)/m$. For $v > 0$, $|v|v = v^2$, so $\partial f/\partial v = -2c_vv_0/m$, and $\partial f/\partial T = 1/m$.

\textbf{Step 3 — the linear model.} With $\delta v = v - v_0$ and $\delta T = T - T_0$: $\;\delta\dot v = -\dfrac{2c_vv_0}{m}\,\delta v + \dfrac{1}{m}\,\delta T$. The drag acts like a linear damper of strength $2c_vv_0$.

\textbf{Step 4 — numbers.} In the full-thrust climb of \lessonref{les:windup}, $v_0 = \SI{2.96}{m/s}$ and $c_v = 0.25$: a damping of \SI{1.48}{N.s/m} and a time constant of $m/(2c_vv_0) = \SI{0.68}{s}$.

\textbf{Step 5 — at hover.} $v_0 = 0$: the derivative $2c_v|v|$ vanishes. To first order there is no drag at all. This is why the linear models of Part~I leave it out — and why the linearisation of the P-only loop cannot see that the drag will, in the end, bring the drone to rest (Example~\ref{ex:G-algebraic}).
\end{example}

\paragraph{How far is \enquote{near}?} A linearisation comes without a range of validity; that must be checked for each use. The horizontal axis of the quadrotor replaces $g\tan\theta$ by $g\,\theta$ (\cref{thm:tilt-chain}). At $10^\circ$ the error is 1\,\%, at $35^\circ$, the tilt limit, 15\,\%: tolerable for feedback, which corrects it, and too much for a feedforward that should be exact.

\section{What stability means}

The definitions are those of \lessonref{les:spring}: an equilibrium is \emph{stable} if solutions that start close stay close, \emph{asymptotically stable} if they also converge to it, \emph{globally} so if they do from every initial state. Two more notions are needed.

\begin{definition}{Exponential stability, region of attraction}
An equilibrium $\vx^*$ is \term{exponentially stable} if there are $c, \alpha > 0$ such that $\|\vx(t) - \vx^*\| \le c\,e^{-\alpha t}\|\vx(0) - \vx^*\|$ for all initial states near it. Its \term{region of attraction} is the set of all initial states from which the solution converges to it.
\end{definition}

For linear systems asymptotic and exponential stability are the same thing, and the region of attraction is everything (\cref{thm:spring-linear}). For nonlinear systems neither is true, and both differences matter in practice: convergence can be slower than any exponential, and it can be lost by starting too far away.

\section{Lyapunov functions}

Aleksandr Lyapunov's idea of 1892 is to avoid solving the equation. Find a function of the state that is positive everywhere except at the equilibrium and that decreases along every solution. A quantity that is bounded below and always decreasing must come to rest, and if it can rest only at the equilibrium, the state must go there.

Take the equilibrium at the origin. For a function $V(\vx)$ the rate of change along a solution follows from the chain rule, without knowing the solution:
\begin{equation}\label{eq:G-vdot}
  \dot V(\vx) = \frac{\partial V}{\partial\vx}\,f(\vx) = \sum_i \frac{\partial V}{\partial x_i}\,f_i(\vx) .
\end{equation}

\begin{theorem}[Lyapunov's direct method, with a rate]
\label{thm:G-lyap}
Let $V$ be continuously differentiable with $V(0) = 0$.
\begin{enumerate}[label=(\roman*)]
  \item If $V(\vx) > 0$ and $\dot V(\vx) \le 0$ for all $\vx \ne 0$ near the origin, the origin is stable.
  \item If moreover $\dot V(\vx) < 0$ for all $\vx \ne 0$ near the origin, it is asymptotically stable.
  \item If there are positive constants with $c_1\|\vx\|^2 \le V(\vx) \le c_2\|\vx\|^2$ and $\dot V(\vx) \le -c_3\|\vx\|^2$, it is exponentially stable, and $\|\vx(t)\| \le \sqrt{c_2/c_1}\;e^{-(c_3/2c_2)\,t}\,\|\vx(0)\|$.
  \item If the conditions hold for all $\vx$ and $V(\vx) \to \infty$ as $\|\vx\| \to \infty$, the conclusions hold globally.
\end{enumerate}
(Khalil, 2002, ch.~4.)
\end{theorem}

The rate in (iii) takes two lines and shows how such proofs work: $\dot V \le -c_3\|\vx\|^2 \le -(c_3/c_2)V$, so $V(t) \le V(0)\,e^{-(c_3/c_2)t}$; then $c_1\|\vx\|^2 \le V$ on the left and $V(0) \le c_2\|\vx(0)\|^2$ on the right, and a square root.

A function that satisfies (i) is a \term{Lyapunov function}. The theorem does not say how to find one. For mechanical systems the energy is the first candidate (\lessonref{les:damper}). For linear systems there is a recipe.

\subsection{Quadratic forms}
The simplest positive functions of a vector are quadratic: $V(\vx) = \vx^\T P\vx$ with a symmetric matrix $P$.

\begin{definition}{Positive definite matrix}
A symmetric matrix $P$ is \term{positive definite}, written $P \succ 0$, if $\vx^\T P\vx > 0$ for every $\vx \ne 0$. Equivalently: all its eigenvalues are positive; or all its leading principal minors (the determinants of its upper-left blocks) are positive. For every $\vx$,
\begin{equation}\label{eq:G-rayleigh}
  \lambda_{min}(P)\,\|\vx\|^2 \le \vx^\T P\vx \le \lambda_{max}(P)\,\|\vx\|^2 .
\end{equation}
\end{definition}

The level sets $\vx^\T P\vx = c$ of a positive definite form are ellipses (ellipsoids), with their axes along the eigenvectors of $P$ and half-lengths $\sqrt{c/\lambda_i}$.

\begin{theorem}[The Lyapunov equation]
\label{thm:G-lyapeq}
All eigenvalues of $A$ have negative real parts if and only if for some (and then for every) symmetric $Q \succ 0$ the equation
\begin{equation}\label{eq:G-lyapeq}
  A^\T P + P A = -Q
\end{equation}
has a symmetric solution $P \succ 0$. The solution is then unique, and $V = \vx^\T P\vx$ satisfies $\dot V = -\vx^\T Q\vx$ along $\dot{\vx} = A\vx$. (Khalil, 2002, ch.~4.)
\end{theorem}

The last claim is a one-line calculation: $\dot V = \dot{\vx}^\T P\vx + \vx^\T P\dot{\vx} = \vx^\T(A^\T P + PA)\vx$. So for a stable linear system a Lyapunov function with a strictly negative derivative always exists, and it is found by solving \emph{linear} equations for the entries of $P$.

\begin{example}[A Lyapunov function for the P–D loop]
\label{ex:G-P}
Solve \cref{eq:G-lyapeq} for the default altitude loop, $A = \begin{psmallmatrix}0 & 1\\ -10 & -7\end{psmallmatrix}$, with $Q = \Id$.

\textbf{Step 1 — unknowns.} $P = \begin{psmallmatrix}p_{11} & p_{12}\\ p_{12} & p_{22}\end{psmallmatrix}$: three numbers.

\textbf{Step 2 — multiply out.} $A^\T P = \begin{pmatrix}-10\,p_{12} & -10\,p_{22}\\ p_{11} - 7p_{12} & p_{12} - 7p_{22}\end{pmatrix}$, and $PA$ is its transpose.

\textbf{Step 3 — three equations.} Entry $(1,1)$: $-20\,p_{12} = -1$. Entry $(2,2)$: $2\,(p_{12} - 7p_{22}) = -1$. Entry $(1,2)$: $p_{11} - 7p_{12} - 10\,p_{22} = 0$.

\textbf{Step 4 — solve.} $p_{12} = 0.05$; $p_{22} = (0.05 + 0.5)/7 = 0.0786$; $p_{11} = 0.35 + 0.786 = 1.136$.

\textbf{Step 5 — is it positive definite?} Leading minors: $1.136 > 0$ and $1.136 \cdot 0.0786 - 0.05^2 = 0.0867 > 0$. Yes — as it must be, since the poles are $-2$ and $-5$.

\textbf{Step 6 — the rate it proves.} The eigenvalues of $P$ are $0.0762$ and $1.138$. \Cref{thm:G-lyap}\,(iii) with $c_1 = 0.0762$, $c_2 = 1.138$, $c_3 = 1$ gives $\|\vx(t)\| \le 3.9\,e^{-0.44t}\|\vx(0)\|$. True, and pessimistic: the slowest pole decays as $e^{-2t}$. A Lyapunov function proves stability; the rates it yields are bounds, and other choices of $Q$ give other bounds.
\end{example}

Compare with the energy $V = \tfrac12m\dot x^2 + \tfrac12K_px^2$ of \lessonref{les:damper}. It is also a quadratic form, with $P = \tfrac12\diag(K_p, m)$, but its derivative $-K_d\dot x^2$ is zero on the whole line $\dot x = 0$: case (i) of the theorem, not (ii). The energy is the \emph{natural} Lyapunov function and the $P$ of Example~\ref{ex:G-P} the \emph{sufficient} one. \Cref{fig:G-lyap} shows both.

\simfig{tb_lyap}{Left: a trajectory of the default P–D loop (blue) crossing the level sets of two Lyapunov functions. The ellipses of $\vx^\T P\vx$ from Example~\ref{ex:G-P} (solid) are crossed inwards everywhere. The ellipses of the energy (dashed) are touched tangentially wherever the velocity is zero. Right: the swing of the drone with air drag and no damper. Dots: peaks of a numerical solution of $\ddot x + 0.25|\dot x|\dot x + 10x = 0$; line: the algebraic decay of Example~\ref{ex:G-algebraic}; dashed: an exponential for comparison.}{fig:G-lyap}

\section{When the derivative is only non-positive}

Energy-like functions usually give $\dot V \le 0$ and no more: dissipation acts on velocities, not on positions. Two theorems close the gap.

\subsection{The invariance principle}
\Cref{thm:damper-lasalle}: if $\dot V \le 0$ everywhere and the only solution that can stay for ever in the set where $\dot V = 0$ is the equilibrium itself, the equilibrium is asymptotically stable. The recipe is always the same: assume $\dot V = 0$ for all time, insert that into the equations of motion, and see what is left.

\begin{example}[Stable, but not exponentially]
\label{ex:G-algebraic}
The ideal P-only loop with air drag and no D term: $m\ddot x + c_v|\dot x|\dot x + K_px = 0$. Its linearisation at rest is the undamped oscillator, which decides nothing (Example~\ref{ex:G-drag}). Does the drone come to rest, and how fast?

\textbf{Step 1 — the candidate.} $V = \tfrac12m\dot x^2 + \tfrac12K_px^2$, positive definite and unbounded at infinity.

\textbf{Step 2 — its derivative.} $\dot V = \dot x\,(m\ddot x + K_px) = -c_v|\dot x|^3 \le 0$.

\textbf{Step 3 — invariance.} $\dot V = 0$ means $\dot x = 0$. A solution that keeps $\dot x = 0$ has $\ddot x = 0$, and the equation leaves $K_px = 0$. Only the origin stays. By \cref{thm:damper-lasalle} it is globally asymptotically stable.

\textbf{Step 4 — how fast.} Use the energy balance of Example~\ref{ex:spring-limit}, now without the motor lag: an oscillation of amplitude $A$ loses $\tfrac83c_vA^3\omega^2$ per period, and $E = \tfrac12K_pA^2$. Dividing by the period as in \cref{eq:spring-amp}:
\begin{equation*}
  \dot A = -\kappa\,A^2, \qquad \kappa = \frac{4\,c_v\,\omega}{3\pi\,m} = \frac{4 \cdot 0.25 \cdot 3.16}{3\pi} = \SI{0.336}{m^{-1}s^{-1}} .
\end{equation*}

\textbf{Step 5 — solve.} Separating variables, $-\dd A/A^2 = \kappa\,\dd t$, gives $1/A - 1/A_0 = \kappa t$:
\begin{equation}\label{eq:G-algebraic}
  A(t) = \frac{A_0}{1 + \kappa A_0\,t} .
\end{equation}

\textbf{Step 6 — read it.} From $A_0 = \SI{1}{m}$: \SI{0.23}{m} after ten seconds, \SI{0.029}{m} after a hundred. The decay is like $1/t$: slower than every exponential, because the smaller the swing, the weaker the drag. A numerical solution of the full equation has peaks of $0.1427$ at \SI{17.9}{s} and $0.0297$ at \SI{97.4}{s}; \cref{eq:G-algebraic} gives the same four digits (\cref{fig:G-lyap}, right).
\end{example}

The example is worth remembering as a counterexample to a common belief: asymptotic stability does \emph{not} imply that errors shrink by a fixed percentage per second. Only exponential stability does, and only exponential stability survives small unmodelled effects — here the motor lag, which turns this slowly converging loop into the limit cycle of \lessonref{les:spring}.

\subsection{Barbalat's lemma}
The invariance principle needs a system that does not depend on time explicitly. An adaptive controller, whose estimates change as it flies, or a loop that tracks a moving reference, is not of that kind. There the tool is a lemma from analysis.

\begin{theorem}[Barbalat's lemma, in the form used for Lyapunov arguments]
\label{thm:G-barbalat}
Let $V(t) \ge 0$ be differentiable with $\dot V(t) \le 0$, and let $\dot V$ be uniformly continuous — which holds, in particular, if $\ddot V$ is bounded. Then $\dot V(t) \to 0$ as $t \to \infty$. (Khalil, 2002, ch.~8.)
\end{theorem}

The use is this. One finds $V \ge 0$ with $\dot V = -(\text{something})^2 \le 0$; then $V$ is bounded, hence so are the signals it contains; from that one checks that $\ddot V$ is bounded; and the lemma yields $(\text{something}) \to 0$. Typically the \enquote{something} is the tracking error, while the parameter estimates inside $V$ are only shown to stay bounded — which is exactly the guarantee an adaptive law gives (Lessons~II.18 and~III.17).\pitfall{A function can have a limit while its derivative does not go to zero: $\sin(t^2)/t$ tends to zero and its derivative keeps oscillating with amplitude 2. The boundedness of $\ddot V$ is what excludes this.}

\section{How far does stability reach?}

Linearisation and the theorems above are local: they hold \enquote{near} the equilibrium. For a loop with limits the engineering question is how large a disturbance it survives. Lyapunov functions answer that too.

\begin{theorem}[An estimate of the region of attraction]
\label{thm:G-roa}
Let $V$ be positive definite and let the set $\Omega_c = \{\vx : V(\vx) \le c\}$ be bounded. If $\dot V(\vx) < 0$ for all $\vx \ne 0$ in $\Omega_c$, then every solution that starts in $\Omega_c$ stays in it and converges to the origin. (Khalil, 2002, ch.~8.)
\end{theorem}

The reasoning is simple: $V$ cannot increase, so a solution cannot cross the boundary $V = c$ outwards; inside, $V$ keeps falling until the state is at the origin. The estimate is the largest level set of $V$ that fits into the region where $\dot V$ is negative. It is always conservative, and a better $V$ gives a larger set.

\begin{example}[A vehicle that is unstable on its own]
\label{ex:G-roa}
A rocket balancing on its engine falls over, when uncontrolled, as $\dot x = a\,x$ with $a = \SI{2}{s^{-1}}$ ($x$ is its tilt, in suitable units). The controller commands $u = -k\,x$ with $k = 6$, and the actuator delivers $\sat(u)$ between $-1$ and $1$. From which tilts can it recover?

\textbf{Step 1 — the loop.} $\dot x = a\,x + \sat(-kx)$. For $|x| \le 1/k = 0.167$ it is linear, $\dot x = -(k - a)x = -4x$: stable.

\textbf{Step 2 — a Lyapunov function.} $V = \tfrac12x^2$, $\dot V = x\,\dot x$.

\textbf{Step 3 — beyond the limit.} For $x > 1/k$ the actuator gives $-1$: $\dot x = ax - 1$, negative only while $x < 1/a = 0.5$. So $\dot V < 0$ for $0 < |x| < 0.5$.

\textbf{Step 4 — the region.} $\Omega_c$ with $c = \tfrac12 \cdot 0.5^2$ is the interval $|x| < 0.5$. From there the rocket recovers. At $|x| > 0.5$ even full actuator authority cannot hold it: $\dot x = ax - 1 > 0$. The estimate is exact here.

\textbf{Step 5 — read it.} The region does not depend on the gain $k$ at all. It is set by the ratio of the actuator's limit to the instability of the plant, $u_{max}/a$. No controller can enlarge it: beyond that tilt the vehicle is lost, and the only defence is never to get there. The drone is kinder: left alone it does not run away exponentially, and a saturated thrust still steers it (\lessonref{les:windup}).
\end{example}

\subsection{Disturbances that never stop}
A loop under turbulence never reaches its equilibrium. What it can promise is to stay close. If a Lyapunov function satisfies
\begin{equation}\label{eq:G-iss}
  \dot V \le -\alpha\,V + \beta\,d^2
\end{equation}
for a disturbance of size $d$, then $V$ decreases whenever $V > \beta d^2/\alpha$: the state enters the set $V \le \beta d_{max}^2/\alpha$ and stays there. It is \term{ultimately bounded}, with a bound proportional to the disturbance. A system with this property for every bounded disturbance is called \emph{input-to-state stable}. Exponentially stable linear systems always are; for nonlinear ones it is a separate property to be proved, and the form \eqref{eq:G-iss} is how (Lesson~III.7).

\section{Limit cycles and averaging}

Not every motion of a nonlinear system ends at an equilibrium. A \term{limit cycle} is an isolated periodic solution; it is stable if neighbouring solutions spiral onto it from both sides. \Lessonref{les:spring} found one by a calculation that deserves its name.

\begin{result}[Averaging a weakly damped oscillator]
For $\ddot x + \omega^2x = \varepsilon\,g(x, \dot x)$ with a small $\varepsilon$, write $x = A\cos(\omega t + \varphi)$ with slowly varying $A$. To first order in $\varepsilon$ the amplitude obeys
\begin{equation}
  \dot A = -\frac{\varepsilon}{2\pi\omega}\int_0^{2\pi} g\big(A\cos\theta,\ -A\omega\sin\theta\big)\,\sin\theta\;\dd\theta ,
\end{equation}
the power of the perturbing force averaged over one period and divided by $\omega^2A$. Equilibria of this scalar equation are limit cycles of the oscillator, and they are stable where the right-hand side decreases through zero. (Khalil, 2002, ch.~10.)
\end{result}

\Cref{eq:spring-amp}, with the motor lag pumping and the drag braking, and \cref{eq:G-algebraic}, with the drag alone, are both instances. The method is justified when the amplitude changes little within one period. Its cousin for loops with a hard nonlinearity — a saturation, a relay — is the describing function of Lesson~III.8, which keeps only the fundamental of the nonlinearity's output (Appendix~C).

\section{Two time scales}

When part of a system is much faster than the rest, the fast part can be replaced by its equilibrium. \Cref{thm:cascade-tikhonov} is the exact statement: for $\dot{\vx} = f(\vx, \symbf z)$, $\varepsilon\dot{\symbf z} = g(\vx, \symbf z)$ with an exponentially stable fast subsystem, the solution of the full system stays within $O(\varepsilon)$ of the solution of the reduced one. The lessons used it three times:
\begin{itemize}
  \item the slow integral and the fast P–D loop (\lessonref{les:droop}), $\varepsilon \approx 0.06$;
  \item the outer and inner loops of the cascade (\lessonref{les:cascade}), $\varepsilon \approx 0.3$;
  \item the amplitude of an oscillation and its phase (the averaging above), $\varepsilon \approx 0.05$.
\end{itemize}
and \lessonref{les:inversion} showed what happens at $\varepsilon \approx 0.8$: the theorem is silent, and the cascade is unstable.

The first correction in $\varepsilon$ can often be computed. For the slow pole of \lessonref{les:droop} it is the second form of \cref{eq:B-slow}: $-\tfrac{K_i}{K_p}\big(1 + \tfrac{K_iK_d}{K_p^2}\big) = -0.0845$, against the quasi-static $-0.080$ and the exact $-0.0850$.

\begin{result}[Appendix G at a glance]
\begin{center}\small
\renewcommand{\arraystretch}{1.4}
\begin{tabularx}{\linewidth}{@{}l>{\raggedright\arraybackslash}X@{}}
linearisation & $A = \partial f/\partial\vx$ at the equilibrium; decides unless eigenvalues sit on the imaginary axis\\
Lyapunov function & $V > 0$, $\dot V = (\partial V/\partial\vx)\,f \le 0$: stable; $\dot V < 0$: asymptotically stable\\
exponential rate & $c_1\|\vx\|^2 \le V \le c_2\|\vx\|^2$, $\dot V \le -c_3\|\vx\|^2$: rate $c_3/(2c_2)$\\
linear systems & solve $A^\T P + PA = -Q$; $V = \vx^\T P\vx$\\
$\dot V \le 0$ only & invariance principle: what can stay where $\dot V = 0$? With time dependence: Barbalat\\
region of attraction & the largest bounded level set of $V$ inside $\{\dot V < 0\}$\\
unstable plant, limited actuator & recoverable region of size $u_{max}/a$, whatever the gain\\
persistent disturbance & $\dot V \le -\alpha V + \beta d^2$: ultimately $V \le \beta d^2/\alpha$\\
weak nonlinearity & average over a period: an equation for the amplitude\\
two time scales & replace the fast part by its equilibrium; error of the order of the ratio\\
\end{tabularx}
\end{center}
\end{result}
